
\documentclass[11pt]{article}

\usepackage{amsmath}
\usepackage{natbib}
\usepackage[colorlinks=true,allcolors=blue]{hyperref}
\usepackage{doi}
\usepackage{url}

\begin{document}

\title{Constructive Logic and Truth}

\author{Charles J. Geyer}

\maketitle

This about the mistake in the document about what one can prove constructively.

First we discuss triple negation reduction.
$$
   \forall P (\neg \neg \neg P \to \neg P)
$$
is provable constructively
(for all $P$ not just $P$ that are mere propositions, trivial because
$\neg P$ is a mere proposition even if $P$ is not)
\citep[Exercise~1.11 in][hereinafter referred to as the HoTT book]{hott-book}.
\defcitealias{hott-book}{HoTT book}

This means if we think of the infinite sequence of statements
\begin{gather*}
   P
   \\
   \neg P
   \\
   \neg \neg P
   \\
   \neg \neg \neg P
   \\
   \neg \neg \neg \neg P
   \\
   \vdots
\end{gather*}
all but the first three logically imply one of the first three.

The question now arises whether for any $P$ we can prove something other
than
\begin{gather*}
   P
   \\
   \neg P
   \\
   \neg \neg P
\end{gather*}
And this is where I stumbled, making the error that Roy Cook pointed out.

What is true is that one cannot prove anything else (that something equivalent
to one of these three).  \citet{geyer-survey} does say this.
But then it goes off the rails asserting that one of the three can always
be proved, but that is not the way constructive mathematics works.  You
might not be able to prove any of these.

Roy Cook mentioned the double negation of LEM is constructively
provable (this is Exercise~{1.13} in the \citetalias{hott-book})
\begin{equation} \label{eq:double-negation-lem}
   \forall P (\neg \neg (P + (\neg P)))
\end{equation}
By the part of De Morgan's laws that are provable constructively
\citepalias[pp.~43-44]{hott-book}, this is equivalent to the
law of noncontradiction
\citep[Section~1.1]{bauer-constructive}
\begin{equation} \label{eq:noncontradiction}
   \neg ((\neg P) \times (\neg \neg P))
\end{equation}

So it is not possible to prove anything other than $P$ or $\neg P$
or $\neg\neg P$ or something that implies one of them (but it is also
possible that you cannot prove any of them).

\begin{thebibliography}{}

\bibitem[Bauer(2017)]{bauer-constructive}
Bauer, A. (2017).
\newblock Five stages of accepting constructive mathematics.
\newblock \emph{Bulletin (New Series) of the American Mathematical Society},
    \textbf{54}, 481--498.
\newblock \doi{10.1090/bull/1556}.

\bibitem[Geyer(2026)]{geyer-survey}
Geyer, C.~J. (2026),
\newblock On generalizations of equality, logic, and set theory in
    mathematics of the past century.
\newblock \url{https://github.com/cjgeyer/Survey}.
\newblock Read by the Foundations Interest Group of the Minnesota Center
    for Philosophy of Science October 1, 2026.

\bibitem[The Univalent Foundations Program(2013)]{hott-book}
The Univalent Foundations Program (2013).
\newblock \emph{Homotopy Type Theory: Univalent Foundations of Mathematics}.
\newblock Institute for Advanced Study, Princeton, NJ.
\newblock \url{https://homotopytypetheory.org/book}.

\end{thebibliography}

\end{document}

